\documentclass[10pt]{article}

% Specify the margins and text size.
\setlength{\textwidth}{6.4in}
\setlength{\textheight}{9.5in}
\setlength{\oddsidemargin}{0pt}
\setlength{\evensidemargin}{0pt}
\setlength{\topmargin}{0pt}
\setlength{\hoffset}{.05in}
\setlength{\voffset}{-1in}

\setlength{\parskip}{5pt}
\setlength{\parindent}{0pt}

% Load some fonts and symbol packages
\usepackage{latexsym}
\usepackage{pifont}       % contains 'star' symbol for counterinsurgency handbook title
\usepackage{yfonts} 
\usepackage{amsmath}
\usepackage{amsfonts}

\usepackage{graphicx}     % actually, this is loaded by pstricks
\usepackage[T1]{fontenc}
\usepackage{ifthen}
%\usepackage{pstricks,pst-grad,pst-text,pst-node,multido,pst-plot,calc,pst-3dplot}
\usepackage{pstricks,pst-grad,pst-text,pst-node,multido,pst-plot,pst-func,calc,pst-3dplot}
%\usepackage[all]{xy}
%\usepackage{animate}

% The hyperref package inserts links into the pdf.
\definecolor{MyLinkColor}{rgb}{.1,.2,1}
\definecolor{MyCiteColor}{rgb}{.1,1,.2}
\definecolor{MyURLColor}{rgb}{.4,.4,.4}
\usepackage[backref=true,pagebackref=false,hyperindex,colorlinks=true,
  linkcolor=MyLinkColor,urlcolor=MyURLColor]{hyperref}


% The tweaklist package is something I found on the web.  It provides a simple interface
% for making changes to spacing used in the itemize and enumerate environments.  Comment
% this out if you don't care to use tweaklist.
\usepackage{tweaklist}
\renewcommand{\itemhook}{\setlength{\parskip}{2pt}\setlength{\parsep}%
{1pt}\setlength{\topsep}{0pt}\setlength{\itemsep}{0pt}}

\newcommand{\U}{\underline{\hspace{5pt}}}

\usepackage{listings}
\newcommand{\Z}{\hphantom{0}}

\begin{document}
\pagestyle{empty}
\lstset{language=R, showspaces=false, showstringspaces=false}
\href{http://www.su.edu}{\includegraphics[height=1.75cm]{sulogo.eps}}
\vspace{-1.79cm}

{{\ }\hfill\small
\begin{tabular}{cl}
& Math 207\\
& Introduction to Statistics\\
\end{tabular}
}
\setlength{\baselineskip}{1.05\baselineskip}
\newcommand{\HH}{\hspace{20pt}\hphantom{a) } }


\begin{center}
\textbf{\large  $z$- and $t$-Test Exercises}
\end{center}
\smallskip

1. A coin is tossed 10,000 times, and it lands heads 5,167 times.  Is the chance
of heads equal to 50\%?  Or are there too many heads for that?
State the null and alternative hypotheses.  Compute $z$ and $P$.
\vspace{2in}

2. Use the table below to find the specified 
area under the $t$-distribution (Student's curve).

\hspace{20pt} a) To the left of $-1.0$ if degrees of freedom is $3$.
\bigskip

\hspace{20pt} b) To the right of $1.0$ if degrees of freedom is $3$.
\bigskip

\hspace{20pt} c) To the right of $1.5$ if degrees of freedom is $2$.
\bigskip

\hspace{20pt} d) To the left of $0.5$ if degrees of freedom is $2$.
\bigskip

\begin{tabular}{ccc}
\begin{tabular}{|c|cccc|}\hline
Degrees\vphantom{\LARGE Y} & &&&\\
of      & &\rput(20pt, 10pt){$t$}&&\\
Freedom & $-2.0$ & $-1.5$ & $-1.0$ & $-0.5$\\[3pt]\hline
1       & $0.148$ & $0.187$ & $0.250$ & $0.352$\vphantom{\LARGE Y}\\
2       & $0.092$ & $0.136$ & $0.211$ & $0.333$\\
3       & $0.070$ & $0.115$ & $0.196$ & $0.326$\\\hline
\end{tabular}
&\hspace{1in} &
\definecolor{lg}{rgb}{.8,.8,.8}
\begin{pspicture}(-2.5,0.5)(2.5,1.5)
\psframe[linewidth=0.02](-3.2,-0.88)(3.2,2.05)
\rput(0,1.7){$t$ distribution with $2$ degrees of freedom}
\psset{yunit=4}
\pscustom[linestyle=none,fillstyle=solid,fillcolor=lg]{%
\psTDist*[linewidth=0,linecolor=blue,nue=2]{-2.25}{-0.5}
\psline(-0.5,0.2963)(-0.5,0)(-2.25,0)(-2.25,0.0533)
}
\psTDist[linecolor=blue,nue=2]{-2.25}{2.25}
\psline{<->}(-2.5,0)(2.5,0)
\psline(-0.5,0)(-0.5,-0.05)\rput[t](-0.5,-0.08){\footnotesize$-0.5\hphantom{-}$}
\psline(-2,0)(-2,-0.05)    \rput[t](-2,-0.08){\footnotesize$-2\hphantom{-}$}
\psline(-1,0)(-1,-0.05)   % \rput[tr](-1,-0.08){\footnotesize$-1$}
\psline(0,0)(0,-0.05)      \rput[t]( 0,-0.08){\footnotesize$0$}
\psline(1,0)(1,-0.05)      \rput[t]( 1,-0.08){\footnotesize$1$}
\psline(2,0)(2,-0.05)      \rput[t]( 2,-0.08){\footnotesize$2$}
\psbezier{->}(-6.13,-0.05)(-4,0.3)(-2,0.5)(-1,0.1)
\end{pspicture}
\end{tabular}
\bigskip
\bigskip

3. A spectrophotometer is being calibrated for its daily use.  Two measurements are
taken which give readings of 71 and 84.  The device is well-calibrated if its
average measurement is 70.  Perform a $t$-test to conclude whether or not the 
device is well-calibrated.
\vfill
\eject

{\ }

4.  One large course has 900 students, broken down into sections meeting with
30 students each.  The section meetings are led by teaching assistants.  On the final, the
class average is 63, and the SD is 20.  However, in one section the 
average is only 55.  The TA argues this way:

\begin{quote}
If you took 30 students at random from the class, there is a pretty good chance 
they would average
below 55 on the final.  That's what happened to me --- chance variation.
\end{quote}

Is this a good defense?  Answer yes or no, and explain briefly.
\vspace{3in}

5. If you used a $z$-test to answer the previous question, redo it using a $t$-test.
If you used $t$, redo it using $z$.
\vfill
\eject
\href{http://www.su.edu}{\includegraphics[height=1.75cm]{sulogo.eps}}
\vspace{-1.79cm}

{{\ }\hfill\small
\begin{tabular}{cl}
& Math 207\\
& Introduction to Statistics\\
\end{tabular}
}
\medskip

\begin{center}
\textbf{\large  The Two-Sample $z$-Test to Compare Two Boxes}
\end{center}
\smallskip

1. 
Four hundred draws are made at random with replacement from box $A$.  The
average of these draws is 110 and their SD is 60.
Independently, 100 draws are made at random from box $B$.  Their average is 90
and SD is 40.  Formulate the null and alternative hypotheses.  Calculate the $z$
and $P$ values.
% SE1 = 3.  SE2 = 4.   SE_diff = 5.  
\vspace{1.5in}


2. Repeat the previous exercise if the average of the draws from $A$ was
95.
\vspace{1.5in}

3. Repeat Exercise 1 if the number of draws from $A$ was 100 and the number
of draws from $B$ was 36.
\vspace{1.5in}

4. %(\#3 from page 506)
In 1970, 59\% of college freshmen thought that capital punishment
should be abolished;  by 2005, the percentage had dropped to 35\%.
Is the difference real, or can it be explained by chance? You may assume that
the percentages are based on two independent simple random samples, each of
size 1,000.

\vfill
\eject
{\ }

5. %(\#5 from page 507)
A university takes a simple random sample of 132 male students and 279 females;
41\% of the men and 17\% of the women report working more than 10 hours 
during the survey week.  To find out whether the difference in percentages is 
statistically significant, the investigator starts by computing
$z=(41-17)/0.048$.  Is anything wrong?    If so fix it.  Then do the test.
\vspace{2in}


6. %(\#1 from page 518)
Five hundred draws are made at random with replacement from a box
of numbered tickets;  276 are positive.  Someone tells you that
50\% of the tickets in the box show positive numbers.  Do you believe it?  
Answer yes or no, and explain.
\vspace{2in}

\vfill
\eject

\end{document}

